\documentclass[11pt]{article}

\usepackage[margin=1in]{geometry}
\usepackage[T1]{fontenc}
\usepackage[utf8]{inputenc}
\usepackage{lmodern}
\usepackage{microtype}
\usepackage{booktabs}
\usepackage{tabularx}
\usepackage{array}
\usepackage{graphicx}
\usepackage{xcolor}
\usepackage{enumitem}
\usepackage{hyperref}
\usepackage{parskip}

\hypersetup{
  colorlinks=true,
  linkcolor=blue!50!black,
  urlcolor=blue!50!black
}

\newcolumntype{L}[1]{>{\raggedright\arraybackslash}p{#1}}
\newcolumntype{Y}{>{\raggedright\arraybackslash}X}

\title{\textbf{Experiment Report: SALAD on GardensPoint and the Campus Dataset}\\[4pt]
Implementation, Motivation, Results, and Comparison with Existing Baselines}
\author{CMU-Africa Night-Time VPR Project}
\date{April 24, 2026}

\begin{document}
\maketitle

\section*{Objective}
The goal of this experiment series was to integrate \textbf{SALAD} into the current VPR codebase and evaluate whether it improves retrieval performance on:
\begin{itemize}[leftmargin=*]
  \item the \textbf{GardensPoint} day-to-night benchmark, and
  \item the project's \textbf{strict 1-to-1 campus dataset}.
\end{itemize}

The main question was whether a stronger modern descriptor could improve retrieval quality without changing the rest of the benchmark pipeline.

This experiment was motivated by two observations:
\begin{itemize}[leftmargin=*]
  \item The current strongest baselines, especially \textbf{CosPlace} and \textbf{EigenPlaces}, already retrieve the correct place within the top-$5$ quite often, but still struggle to rank it first consistently.
  \item The recent literature review suggested that \textbf{SALAD} is one of the strongest modern drop-in global descriptor methods and therefore a good first candidate for extension.
\end{itemize}

\section*{What the Original SALAD Paper Tested}
SALAD was introduced in the CVPR 2024 paper \textit{Optimal Transport Aggregation for Visual Place Recognition} by Izquierdo and Civera. In the released setup, the model was:
\begin{itemize}[leftmargin=*]
  \item fine-tuned on \textbf{GSV-Cities},
  \item monitored on \textbf{Pittsburgh30k-test},
  \item and evaluated on several standard VPR benchmarks rather than on GardensPoint alone.
\end{itemize}

The main datasets reported in the paper include:
\begin{itemize}[leftmargin=*]
  \item \textbf{MSLS Challenge}
  \item \textbf{MSLS Val}
  \item \textbf{NordLand}
  \item \textbf{Pittsburgh250k-test}
  \item \textbf{SPED}
  \item \textbf{SF-XL}
\end{itemize}

The strongest published SALAD variant in the paper is the \textbf{DINOv2 SALAD (8192 + 256)} model. The headline reported numbers were:

\begin{center}
\begin{tabular}{L{3.3cm}L{1.7cm}L{1.7cm}}
\toprule
\textbf{Dataset} & \textbf{R@1} & \textbf{R@5} \\
\midrule
MSLS Challenge & 75.0 & 88.8 \\
MSLS Val & 92.2 & 96.4 \\
NordLand & 76.0 & 89.2 \\
Pittsburgh250k-test & 95.1 & 98.5 \\
SPED & 92.1 & 96.2 \\
\bottomrule
\end{tabular}
\end{center}

They also reported:
\begin{itemize}[leftmargin=*]
  \item \textbf{MSLS Challenge}: R@10 = 91.3
  \item \textbf{MSLS Val}: R@10 = 97.0
  \item \textbf{SF-XL Test v1}: R@1 = 88.6
  \item \textbf{SF-XL Test v2}: R@1 = 94.8
\end{itemize}

This context matters for our reproduction: the published SALAD results come from \textbf{large-scale standard VPR benchmarks}, while our current experiments are on \textbf{GardensPoint} and the project's \textbf{custom campus datasets}. So our report should be interpreted as a \textbf{local integration and transfer experiment}, not a direct re-creation of the paper's full benchmark suite.

\section*{Why GardensPoint Is Still the Right Primary Benchmark}
Although the original SALAD paper reports results on larger standard datasets such as \textbf{MSLS} and \textbf{Pittsburgh250k}, \textbf{GardensPoint} remains the most appropriate primary benchmark for this project.

The reason is simple: the main project question is \textbf{day-to-night visual place recognition}. GardensPoint directly matches that setting because it evaluates retrieval across a strong illumination shift between day and night imagery. By contrast, datasets such as \textbf{Pittsburgh} are very useful for general VPR benchmarking, but they are less aligned with the specific nocturnal transfer problem we are trying to solve.

For this project, the benchmark roles are therefore:
\begin{itemize}[leftmargin=*]
  \item \textbf{GardensPoint}: primary controlled benchmark for day-to-night evaluation,
  \item \textbf{campus dataset}: local transfer and deployment-oriented validation,
  \item \textbf{large-scale standard datasets such as Pittsburgh}: optional auxiliary benchmarks for broader descriptor sanity checks, but not the main evidence for the night-time claim.
\end{itemize}

So even though GardensPoint is smaller than some standard VPR datasets, it is still the correct main benchmark because it is the one most directly aligned with the problem definition of this work.

\section*{What Was Implemented}
The SALAD integration was done in a modular way so that it can be reused by both the benchmark scripts and the live VPR pipeline.

\subsection*{Step 1: Shared extractor factory}
Instead of adding SALAD as another one-off code path, the descriptor registry was centralized in:
\begin{itemize}[leftmargin=*]
  \item \texttt{feature\_extraction/factory.py}
\end{itemize}

This file now:
\begin{itemize}[leftmargin=*]
  \item defines the shared list of supported descriptors,
  \item exposes \texttt{create\_feature\_extractor()},
  \item keeps descriptor-specific loading logic in one place.
\end{itemize}

\subsection*{Step 2: Reusable Torch Hub extractor base}
A reusable base class was added in:
\begin{itemize}[leftmargin=*]
  \item \texttt{feature\_extraction/feature\_extractor\_torchhub.py}
\end{itemize}

This wrapper handles:
\begin{itemize}[leftmargin=*]
  \item device selection,
  \item image preprocessing,
  \item small-batch and loader-based feature extraction,
  \item cleaner failure messages for missing dependencies.
\end{itemize}

\subsection*{Step 3: SALAD-specific wrapper}
The SALAD extractor itself was added in:
\begin{itemize}[leftmargin=*]
  \item \texttt{feature\_extraction/feature\_extractor\_salad.py}
\end{itemize}

The wrapper uses the official Torch Hub entrypoint and loads the released \texttt{dinov2\_salad} checkpoint. In the current setup this produces an \textbf{8448-dimensional global descriptor}.

\subsection*{Step 4: Unified benchmark integration}
The benchmark scripts were updated to use the shared extractor factory:
\begin{itemize}[leftmargin=*]
  \item \texttt{demo.py}
  \item \texttt{test\_campus\_dataset.py}
  \item \texttt{test\_campus\_dataset\_place\_level.py}
\end{itemize}

This means SALAD is now selectable in:
\begin{itemize}[leftmargin=*]
  \item GardensPoint,
  \item St Lucia,
  \item SFU,
  \item the strict 1-to-1 campus evaluation,
  \item the place-level campus evaluation.
\end{itemize}

\subsection*{Step 5: Practical runtime fix on Apple Silicon}
During testing, SALAD failed on \texttt{MPS} because the DINOv2 backbone uses a bicubic upsampling path that is not implemented on MPS in the current PyTorch release.

To make the experiment usable immediately, SALAD was patched to:
\begin{itemize}[leftmargin=*]
  \item \textbf{fall back to CPU automatically on MPS devices},
  \item while leaving the other PyTorch extractors free to keep using MPS.
\end{itemize}

This made the experiment slower, but reliable.

\section*{How the Experiments Were Run}
The first benchmark run used the standard repo entrypoint:

\begin{verbatim}
python demo.py --descriptor SALAD --dataset GardensPoint --save_results
\end{verbatim}

The steps executed by the script were:
\begin{enumerate}[leftmargin=*]
  \item load the GardensPoint day-right / night-right benchmark,
  \item compute reference-set descriptors using SALAD,
  \item compute query-set descriptors using SALAD,
  \item normalize the descriptors,
  \item compute cosine similarity matrix $S$,
  \item apply the existing matching and evaluation code,
  \item compute AUC, R@100P, and Recall@K.
\end{enumerate}

The second run used the strict campus benchmark entrypoint:

\begin{verbatim}
python test_campus_dataset.py --descriptor SALAD --save_results
\end{verbatim}

This followed the same descriptor-extraction and cosine-similarity flow, but evaluated against the custom campus ground-truth labels where only \textbf{33 queries} have exact matches and \textbf{31 queries} are marked as no-match.

\section*{Observed SALAD Result on GardensPoint}
The run produced the following result summary:

\begin{center}
\begin{tabular}{L{2.8cm}L{2.2cm}}
\toprule
\textbf{Metric} & \textbf{SALAD Result} \\
\midrule
AUC & 0.924 \\
R@100P & 0.385 \\
R@1 & 0.450 \\
R@5 & 0.955 \\
R@10 & 0.990 \\
\bottomrule
\end{tabular}
\end{center}

In addition, descriptor extraction was noticeably slower than the current ResNet-based baselines because the run used \textbf{CPU fallback}:
\begin{itemize}[leftmargin=*]
  \item database descriptors: about \textbf{3 min 19 s}
  \item query descriptors: about \textbf{3 min 08 s}
\end{itemize}

\begin{figure}[h]
\centering
\includegraphics[width=\linewidth,height=8.5cm,keepaspectratio]{../../output_images/runs/benchmarks/0010_demo_gardenspoint_salad_20260424_161122/matches_examples.png}
\caption{Example qualitative result from the exact saved SALAD run on GardensPoint.}
\end{figure}

\section*{Strict Review of the GardensPoint Result}
Because the GardensPoint SALAD result is much stronger than the earlier baselines, it deserves a stricter review before being trusted.

\subsection*{What looks correct}
After reviewing the code path, there is \textbf{no obvious evidence of data leakage or a broken evaluation loop}:
\begin{itemize}[leftmargin=*]
  \item the dataset loader reads \texttt{day\_right} and \texttt{night\_right} from separate folders,
  \item the model only sees image tensors, not filenames,
  \item the similarity matrix is computed in the standard way by cosine similarity after descriptor normalization,
  \item and \textbf{Recall@K} is evaluated against the exact hard ground truth.
\end{itemize}

So the result does \textbf{not} look like a trivial coding mistake such as mixing queries and references, using labels during inference, or directly comparing identical images.

\subsection*{What should be treated cautiously}
Even so, a strict reviewer should \textbf{not} interpret all reported metrics equally.

\begin{itemize}[leftmargin=*]
  \item \textbf{GardensPoint uses identity hard ground truth and a wide soft ground truth band.} In the loader, \texttt{GThard} is the identity matrix, and \texttt{GTsoft} is created by a convolution with a \texttt{17x1} kernel, which effectively allows about \textbf{$\pm 8$ nearby reference frames} as soft matches.
  \item \textbf{AUC and R@100P are therefore relatively lenient.} In \texttt{evaluation/metrics.py}, entries that are true only in \texttt{GTsoft} are pushed down to \texttt{S.min()} and ignored during threshold evaluation. That makes the precision-recall metrics more forgiving than the exact-match top-$k$ metrics.
  \item \textbf{The most trustworthy summary numbers are therefore R@1, R@5, and R@10.} These use the exact hard ground truth rather than the softened threshold logic.
  \item \textbf{GardensPoint is still a small single-route benchmark.} So even if the result is valid, it should be treated as evidence of strong controlled day-to-night retrieval, not as proof of broad real-world robustness.
\end{itemize}

\subsection*{Professor-style conclusion}
From a strict computer-vision review perspective, the GardensPoint SALAD result looks \textbf{plausible}, not impossible:
\begin{itemize}[leftmargin=*]
  \item \textbf{R@1 = 0.450} is strong, but not suspiciously close to perfect,
  \item \textbf{R@5 = 0.955} and \textbf{R@10 = 0.990} fit the existing project narrative that the correct place is often already in the top few candidates,
  \item the most inflated-looking number is \textbf{AUC = 0.924}, but that can be explained by the wide soft ground-truth treatment rather than by a clear coding bug.
\end{itemize}

So the appropriate trust statement is:
\begin{quote}
The SALAD GardensPoint result appears to be a \textbf{valid and strong benchmark result}, but the most defensible evidence is the \textbf{Recall@K} improvement, not the threshold-based metrics alone.
\end{quote}

One additional reproducibility note is important: this report now references the \textbf{run-specific saved artifacts} from \texttt{output\_images/runs/...} rather than the floating top-level files in \texttt{output\_images/}. That avoids a common reporting mistake where later runs silently overwrite earlier qualitative figures.

\section*{Observed SALAD Result on the Strict Campus Dataset}
The strict campus run produced the following result summary:

\begin{center}
\begin{tabular}{L{2.8cm}L{2.2cm}}
\toprule
\textbf{Metric} & \textbf{SALAD Result} \\
\midrule
AUC & 0.504 \\
R@100P & 0.030 \\
R@1 & 0.818 \\
R@5 & 1.000 \\
R@10 & 1.000 \\
TP (thresholded) & 7 \\
FP (thresholded) & 6 \\
\bottomrule
\end{tabular}
\end{center}

The qualitative match visualization selected:
\begin{itemize}[leftmargin=*]
  \item 3 correct examples from 7 thresholded true-positive pairs,
  \item 5 wrong examples from 6 thresholded false-positive pairs.
\end{itemize}

The no-match analysis also showed that several \texttt{-npm} queries still had relatively high-similarity top matches, for example:
\begin{itemize}[leftmargin=*]
  \item Query 0 $\rightarrow$ DB 7 with similarity 0.700
  \item Query 1 $\rightarrow$ DB 7 with similarity 0.683
  \item Query 4 $\rightarrow$ DB 10 with similarity 0.627
\end{itemize}

This indicates that SALAD is strong at \textbf{retrieval}, but that open-set rejection remains difficult under the strict campus protocol.

\begin{figure}[h]
\centering
\includegraphics[width=\linewidth,height=8.5cm,keepaspectratio]{../../output_images/runs/campus_strict/0004_campus_strict_salad_20260424_150007/matches_examples.png}
\caption{Example qualitative result from the exact saved SALAD run on the strict campus dataset.}
\end{figure}

\section*{Observed SALAD Result on the Place-Level Campus Dataset}
The place-level campus run produced the following result summary:

\begin{center}
\begin{tabular}{L{2.8cm}L{2.2cm}}
\toprule
\textbf{Metric} & \textbf{SALAD Result} \\
\midrule
AUC & 0.653 \\
R@100P & 0.130 \\
R@1 & 0.969 \\
R@5 & 1.000 \\
R@10 & 1.000 \\
TP (thresholded) & 23 \\
FP (thresholded) & 0 \\
\bottomrule
\end{tabular}
\end{center}

This run used the relabeled place-level dataset, where:
\begin{itemize}[leftmargin=*]
  \item all \textbf{64 queries} have valid place-level matches,
  \item the original \texttt{-npm} cases have been manually reviewed,
  \item \textbf{12 queries} were rematched from the original no-match set,
  \item and the campus route is organized into \textbf{10 place IDs}.
\end{itemize}

The thresholded behavior is much cleaner here:
\begin{itemize}[leftmargin=*]
  \item \textbf{23 true positives},
  \item \textbf{0 false positives}.
\end{itemize}

Runtime was again dominated by CPU descriptor extraction:
\begin{itemize}[leftmargin=*]
  \item database descriptors: about \textbf{1 min 45 s}
  \item query descriptors: about \textbf{1 min 54 s}
\end{itemize}

The run also showed that \textbf{xFormers was not available}, but the model still executed correctly. So this result reflects a valid SALAD run, just not the most optimized one.

\begin{figure}[h]
\centering
\includegraphics[width=\linewidth,height=8.5cm,keepaspectratio]{../../output_images/runs/campus_place_level/0004_campus_place_level_salad_20260424_151858/matches_examples.png}
\caption{Example qualitative result from the exact saved SALAD run on the place-level campus dataset.}
\end{figure}

\section*{Comparison with Existing Baselines}
The project's earlier reported GardensPoint baselines were:

\begin{center}
\begin{tabular}{L{3.0cm}L{1.2cm}L{1.2cm}L{1.2cm}L{1.2cm}}
\toprule
\textbf{Descriptor} & \textbf{AUC} & \textbf{R@1} & \textbf{R@5} & \textbf{R@10} \\
\midrule
CosPlace & 0.592 & 0.350 & 0.840 & 0.925 \\
EigenPlaces & 0.676 & 0.370 & 0.840 & 0.940 \\
AlexNet conv3 & 0.195 & 0.300 & 0.665 & 0.745 \\
SAD & 0.030 & 0.090 & 0.225 & 0.305 \\
SALAD & \textbf{0.924} & \textbf{0.450} & \textbf{0.955} & \textbf{0.990} \\
\bottomrule
\end{tabular}
\end{center}

Compared with the strongest prior baseline (\textbf{EigenPlaces}), SALAD improved:

\begin{center}
\begin{tabular}{L{2.8cm}L{2.0cm}}
\toprule
\textbf{Metric} & \textbf{Change vs EigenPlaces} \\
\midrule
AUC & +0.248 \\
R@1 & +0.080 \\
R@5 & +0.115 \\
R@10 & +0.050 \\
\bottomrule
\end{tabular}
\end{center}

Compared with \textbf{CosPlace}, SALAD improved:

\begin{center}
\begin{tabular}{L{2.8cm}L{2.0cm}}
\toprule
\textbf{Metric} & \textbf{Change vs CosPlace} \\
\midrule
AUC & +0.332 \\
R@1 & +0.100 \\
R@5 & +0.115 \\
R@10 & +0.065 \\
\bottomrule
\end{tabular}
\end{center}

\section*{Comparison with Existing Campus Baselines}
The previously reported strict campus baselines were:

\begin{center}
\begin{tabular}{L{2.8cm}L{1.1cm}L{1.1cm}L{1.0cm}L{1.0cm}L{1.0cm}}
\toprule
\textbf{Descriptor} & \textbf{AUC} & \shortstack{\textbf{R@}\\\textbf{100P}} & \textbf{R@1} & \textbf{R@5} & \shortstack{\textbf{R@}\\\textbf{10}} \\
\midrule
CosPlace & 0.470 & 0.030 & 0.727 & 1.000 & 1.000 \\
EigenPlaces & 0.443 & 0.000 & 0.788 & 1.000 & 1.000 \\
NetVLAD & 0.242 & 0.000 & 0.394 & 0.879 & 0.939 \\
AlexNet & \textbf{0.512} & \textbf{0.061} & 0.758 & 0.939 & 0.939 \\
HDC-DELF & 0.488 & 0.030 & 0.788 & 0.939 & 0.970 \\
SAD & 0.106 & 0.000 & 0.212 & 0.545 & 0.727 \\
SALAD & 0.504 & 0.030 & \textbf{0.818} & \textbf{1.000} & \textbf{1.000} \\
\bottomrule
\end{tabular}
\end{center}

Compared with the strict campus baselines, SALAD shows a mixed but useful pattern:
\begin{itemize}[leftmargin=*]
  \item it has the \textbf{best R@1} among the strict campus descriptors tested so far,
  \item it ties the strongest methods at \textbf{R@5} and \textbf{R@10},
  \item it has \textbf{very competitive AUC}, although \textbf{AlexNet} remains slightly higher at 0.512,
  \item it does \textbf{not} improve open-set rejection under the strict protocol, since \textbf{R@100P} remains 0.030.
\end{itemize}

\section*{Comparison with Existing Place-Level Campus Baselines}
The previously reported place-level campus baselines were:

\begin{center}
\begin{tabular}{L{2.8cm}L{1.1cm}L{1.1cm}L{1.0cm}L{1.0cm}L{1.0cm}}
\toprule
\textbf{Descriptor} & \textbf{AUC} & \shortstack{\textbf{R@}\\\textbf{100P}} & \textbf{R@1} & \textbf{R@5} & \shortstack{\textbf{R@}\\\textbf{10}} \\
\midrule
CosPlace & 0.627 & 0.114 & \textbf{0.969} & 1.000 & 1.000 \\
EigenPlaces & \textbf{0.663} & 0.110 & 0.922 & 1.000 & 1.000 \\
NetVLAD & 0.512 & 0.024 & 0.891 & 0.938 & 0.969 \\
PatchNetVLAD & 0.641 & 0.074 & 0.938 & 1.000 & 1.000 \\
AlexNet & 0.387 & 0.074 & 0.859 & 0.969 & 1.000 \\
HDC-DELF & 0.507 & \textbf{0.132} & 0.922 & 0.969 & 0.984 \\
SAD & 0.170 & 0.000 & 0.297 & 0.625 & 0.766 \\
SALAD & 0.653 & 0.130 & \textbf{0.969} & 1.000 & 1.000 \\
\bottomrule
\end{tabular}
\end{center}

Compared with the place-level baselines, SALAD shows a stronger and cleaner picture than on the strict campus setup:
\begin{itemize}[leftmargin=*]
  \item it \textbf{ties the best R@1} at \textbf{0.969} with CosPlace,
  \item it matches the strongest methods at \textbf{R@5} and \textbf{R@10},
  \item it has the \textbf{second-best AUC} at \textbf{0.653}, just below EigenPlaces at 0.663,
  \item it has the \textbf{second-best R@100P} at \textbf{0.130}, just below HDC-DELF at 0.132,
  \item and unlike the strict run, it produced \textbf{0 thresholded false positives}.
\end{itemize}

\section*{Interpretation of the Result}
\subsection*{1. SALAD is clearly stronger than the current baselines on GardensPoint}
The most immediate conclusion is that \textbf{SALAD substantially outperformed all existing reported baselines} on GardensPoint. The AUC jump is especially large, which suggests SALAD is not only retrieving correct places better, but doing so with a stronger precision-recall tradeoff overall.

\subsection*{2. The ranking problem improved, but did not disappear}
The current project framing has emphasized that the main issue is often \textbf{ranking} rather than raw retrieval. SALAD supports that interpretation:
\begin{itemize}[leftmargin=*]
  \item it strongly improves \textbf{R@5} and \textbf{R@10},
  \item and it also improves \textbf{R@1},
  \item but \textbf{R@1 = 0.450} still leaves room for better first-rank consistency.
\end{itemize}

This means the earlier plan around \textbf{sequence matching} still makes sense. SALAD appears to provide a stronger descriptor backbone, but temporal re-ranking may still improve final localization stability.

\subsection*{3. The model is promising, but currently expensive on this machine}
The main practical downside is runtime. On this Apple Silicon setup, SALAD had to run on \textbf{CPU} because of the MPS limitation in the DINOv2 interpolation path. As a result:
\begin{itemize}[leftmargin=*]
  \item the benchmark run was much slower than the current ResNet-based baselines,
  \item and live deployment with SALAD on this machine would likely need either lower processing cadence or additional optimization.
\end{itemize}

So the current interpretation is:
\begin{itemize}[leftmargin=*]
  \item \textbf{excellent benchmark candidate},
  \item \textbf{strong retrieval upgrade},
  \item but \textbf{not yet a cheap live-deployment choice on MPS hardware}.
\end{itemize}

\subsection*{4. The campus result reinforces the retrieval-vs-rejection split}
The strict campus result is especially informative because it separates two behaviors:
\begin{itemize}[leftmargin=*]
  \item \textbf{retrieval quality}: SALAD is very strong, as shown by \textbf{R@1 = 0.818} and perfect \textbf{R@5 / R@10},
  \item \textbf{rejection quality}: SALAD is still limited under strict open-set handling, as shown by \textbf{R@100P = 0.030}.
\end{itemize}

This means SALAD is likely to be a strong backbone for the \textbf{place-level campus evaluation}, but the strict 1-to-1 campus result also warns that descriptor strength alone may not solve the no-match problem.

\subsection*{5. The place-level campus result confirms that SALAD is especially strong when the evaluation credits valid alternate views}
The place-level result is arguably the most balanced campus result for SALAD:
\begin{itemize}[leftmargin=*]
  \item \textbf{R@1 = 0.969} shows that SALAD can rank the correct place first almost all the time once same-place alternate views are credited fairly,
  \item \textbf{AUC = 0.653} shows a very strong overall ranking profile,
  \item \textbf{R@100P = 0.130} shows that precision-recall behavior also improves substantially compared with the strict campus run,
  \item and \textbf{FP = 0} under thresholded place-level matching shows cleaner match decisions.
\end{itemize}

This supports the broader conclusion you already reached with the campus relabeling work: some of the earlier campus errors were caused by the \textbf{evaluation protocol}, not only by descriptor weakness.

\section*{Why This Experiment Matters}
This experiment is important for the project in two ways.

First, it validates the recent research review: \textbf{SALAD was a sensible first method to try}, and the result supports that choice strongly.

Second, it sharpens the next design question. The project no longer only asks:
\begin{quote}
Can we beat the current baselines?
\end{quote}

It now asks:
\begin{quote}
Can we combine a stronger descriptor backbone such as SALAD with the planned lightweight improvements (CLAHE, sequence matching, and live search optimization) without losing deployment practicality?
\end{quote}

\section*{Recommended Next Steps}
Based on these first SALAD experiments, the next steps should be:
\begin{enumerate}[leftmargin=*]
  \item \textbf{Compare strict and place-level SALAD results more explicitly}
  \begin{itemize}[leftmargin=*]
    \item this run is now done for both campus protocols,
    \item the next useful step is to summarize the strict vs place-level gap in slides and discussion.
  \end{itemize}

  \item \textbf{Test SALAD with the planned sequence matching layer}
  \begin{itemize}[leftmargin=*]
    \item this is the most direct way to address the remaining ranking problem.
  \end{itemize}

  \item \textbf{Compare SALAD against the next research candidates}
  \begin{itemize}[leftmargin=*]
    \item especially \textbf{BoQ}, \textbf{AnyLoc}, and \textbf{MixVPR}.
  \end{itemize}

  \item \textbf{Profile live feasibility}
  \begin{itemize}[leftmargin=*]
    \item test lower \texttt{process\_fps},
    \item test ANN search with large descriptors,
    \item decide whether SALAD belongs in the live demo path or mainly in the benchmark path.
  \end{itemize}
\end{enumerate}

\section*{Conclusion}
The SALAD experiments on GardensPoint and the strict campus dataset were successful overall. On GardensPoint, SALAD delivered a \textbf{clear improvement over the current project baselines}, especially in AUC and top-$k$ recall. On the strict campus dataset, SALAD achieved the \textbf{best R@1} seen so far and remained highly competitive in AUC.

The place-level campus result strengthens the story further. Under the fairer multi-view evaluation, SALAD reached:
\begin{itemize}[leftmargin=*]
  \item \textbf{AUC = 0.653}
  \item \textbf{R@1 = 0.969}
  \item \textbf{R@100P = 0.130}
  \item \textbf{0 thresholded false positives}
\end{itemize}

That makes SALAD one of the strongest overall campus descriptors tested so far, especially once same-place alternate views are credited properly.

At the same time, the experiment exposed a practical deployment cost: on the current Apple Silicon setup, SALAD requires \textbf{CPU fallback} and therefore runs much more slowly than the current lighter baselines.

So the correct conclusion is:
\begin{itemize}[leftmargin=*]
  \item \textbf{SALAD is a strong benchmark and campus-retrieval backbone,}
  \item \textbf{it is especially convincing under place-level campus evaluation,}
  \item \textbf{it still does not solve strict no-match rejection by itself,}
  \item and \textbf{live deployment feasibility still needs separate validation.}
\end{itemize}

\end{document}
